% ============================================================
% Question Bank: ch08-01 Populations and Samples (Modified — OK/TX)
% Topic Title: Populations and Samples (Experimental Design Emphasis)
% CCSS: 7.SP.A.1 (OK/TX: experimental design emphasis)
% Questions: 30 (18 MC + 7 SA + 5 Graphical)
% ============================================================

% --- Multiple Choice Questions (18) ---

\begin{practiceQuestion}{m-8-1-q01}{mc}
\begin{questionText}
The first step in designing an experiment is to:
\end{questionText}
\choiceA{Collect data}
\choiceB{Choose a sampling method}
\choiceC{Ask a clear question}
\choiceD{Analyze the results}
\correctAnswer{C}
\explanation{Before anything else, you need to know what you want to find out — start with a clear question.}
\end{practiceQuestion}

\begin{practiceQuestion}{m-8-1-q02}{mc}
\begin{questionText}
A school has $600$ students. A researcher randomly selects $60$ student IDs using a random number generator. This is an example of:
\end{questionText}
\choiceA{A convenience sample}
\choiceB{A voluntary response sample}
\choiceC{A random sample}
\choiceD{A biased sample}
\correctAnswer{C}
\explanation{Using a random number generator gives every student an equal chance of being selected, making it a random sample.}
\end{practiceQuestion}

\begin{practiceQuestion}{m-8-1-q03}{mc}
\begin{questionText}
Which of these is a good statistical question for a school survey?
\end{questionText}
\choiceA{How many students are in the school?}
\choiceB{What is the tallest building in the city?}
\choiceC{How much time do students at our school spend on homework each night?}
\choiceD{What is $2 + 3$?}
\correctAnswer{C}
\explanation{A good statistical question expects variability in the answers and can be answered by collecting data from a sample.}
\end{practiceQuestion}

\begin{practiceQuestion}{m-8-1-q04}{mc}
\begin{questionText}
A teacher wants to know the favorite lunch of all $500$ students. She surveys the first $30$ students who enter the cafeteria. What is the problem?
\end{questionText}
\choiceA{The sample is too large}
\choiceB{The sample is biased — early arrivals may not represent all students}
\choiceC{There is no problem}
\choiceD{Thirty students is a good random sample}
\correctAnswer{B}
\explanation{Surveying only early arrivals is a convenience sample. Students who arrive first may have different preferences than others.}
\end{practiceQuestion}

\begin{practiceQuestion}{m-8-1-q05}{mc}
\begin{questionText}
Which of the following is NOT one of the five steps of designing an experiment?
\end{questionText}
\choiceA{Ask a question}
\choiceB{Identify the population}
\choiceC{Make up the data you expect}
\choiceD{Collect data}
\correctAnswer{C}
\explanation{The five steps are: ask a question, identify the population, choose a sampling method, collect data, analyze and conclude. You never make up data.}
\end{practiceQuestion}

\begin{practiceQuestion}{m-8-1-q06}{mc}
\begin{questionText}
A city wants to know residents' opinions about a new park. Which sampling plan avoids bias?
\end{questionText}
\choiceA{Survey people who already use the current park}
\choiceB{Randomly select $200$ addresses from the city directory}
\choiceC{Ask people at the city council meeting}
\choiceD{Post the survey online and wait for responses}
\correctAnswer{B}
\explanation{Randomly selecting addresses from the full city directory gives every resident an equal chance of being surveyed.}
\end{practiceQuestion}

\begin{practiceQuestion}{m-8-1-q07}{mc}
\begin{questionText}
A company surveys only its own employees about a new product. The population of interest is all potential customers. This sample is:
\end{questionText}
\choiceA{Random and representative}
\choiceB{Biased because employees may have different opinions than customers}
\choiceC{Large enough to be valid}
\choiceD{The best available option}
\correctAnswer{B}
\explanation{Employees may feel loyalty or pressure, which makes their opinions different from the general customer population.}
\end{practiceQuestion}

\begin{practiceQuestion}{m-8-1-q08}{mc}
\begin{questionText}
Describe the population and sample: A factory tests $50$ light bulbs from a batch of $5{,}000$.
\end{questionText}
\choiceA{Population: $50$ bulbs; Sample: $5{,}000$ bulbs}
\choiceB{Population: $5{,}000$ bulbs; Sample: $50$ bulbs}
\choiceC{Population: all light bulbs ever made; Sample: the factory}
\choiceD{Population: the factory; Sample: $50$ bulbs}
\correctAnswer{B}
\explanation{The population is the entire batch ($5{,}000$ bulbs). The sample is the $50$ bulbs tested.}
\end{practiceQuestion}

\begin{practiceQuestion}{m-8-1-q09}{mc}
\begin{questionText}
A student wants to find out which school club is most popular. She asks $10$ students in the science club. Why is this a poor experimental design?
\end{questionText}
\choiceA{The sample size is fine}
\choiceB{Science club members are likely biased toward the science club}
\choiceC{She should ask fewer students}
\choiceD{There is no problem with this design}
\correctAnswer{B}
\explanation{Asking only science club members introduces bias — they are more likely to say the science club is the most popular.}
\end{practiceQuestion}

\begin{practiceQuestion}{m-8-1-q10}{mc}
\begin{questionText}
After collecting data from a random sample, the next step is to:
\end{questionText}
\choiceA{Start over with a new question}
\choiceB{Change the sample to get better results}
\choiceC{Analyze the data and draw conclusions}
\choiceD{Ignore the results}
\correctAnswer{C}
\explanation{After collecting data, you analyze it (find means, percents, patterns) and draw conclusions about the population.}
\end{practiceQuestion}

\begin{practiceQuestion}{m-8-1-q11}{mc}
\begin{questionText}
A random sample of $80$ out of $2{,}000$ students found that $20$ students bicycle to school. Predict how many in the school bicycle to school.
\end{questionText}
\choiceA{$200$}
\choiceB{$400$}
\choiceC{$500$}
\choiceD{$1{,}000$}
\correctAnswer{C}
\explanation{$\frac{20}{80} = 25\%$. Then $25\%$ of $2{,}000 = 500$.}
\end{practiceQuestion}

\begin{practiceQuestion}{m-8-1-q12}{mc}
\begin{questionText}
Why is a small random sample usually better than a large biased sample?
\end{questionText}
\choiceA{Smaller samples are always more accurate}
\choiceB{A random sample is representative even if small; a biased sample misrepresents the population regardless of size}
\choiceC{Bigger samples are always worse}
\choiceD{Bias improves with sample size}
\correctAnswer{B}
\explanation{A random sample, even if small, fairly represents the population. A biased sample systematically misrepresents it, no matter how large.}
\end{practiceQuestion}

\begin{practiceQuestion}{m-8-1-q13}{mc}
\begin{questionText}
A researcher designs an experiment with these steps: (1) question, (2) population, (3) sampling method, (4) data collection, (5) analysis. She wants to know if students prefer chocolate or vanilla. Which step involves tallying the results and finding percents?
\end{questionText}
\choiceA{Step 1}
\choiceB{Step 3}
\choiceC{Step 4}
\choiceD{Step 5}
\correctAnswer{D}
\explanation{Step 5 (analyze and conclude) is where you organize results, calculate statistics, and draw conclusions.}
\end{practiceQuestion}

\begin{practiceQuestion}{m-8-1-q14}{mc}
\begin{questionText}
A survey of $100$ randomly selected adults in a town found that $45\%$ support a new library. If the town has $10{,}000$ adults, about how many support the library?
\end{questionText}
\choiceA{$450$}
\choiceB{$4{,}500$}
\choiceC{$5{,}500$}
\choiceD{$10{,}000$}
\correctAnswer{B}
\explanation{$45\%$ of $10{,}000 = 0.45 \times 10{,}000 = 4{,}500$.}
\end{practiceQuestion}

\begin{practiceQuestion}{m-8-1-q15}{mc}
\begin{questionText}
Which of the following describes a systematic sampling method?
\end{questionText}
\choiceA{Asking every $5$th person on a class list}
\choiceB{Asking only students who raise their hand}
\choiceC{Asking students you happen to see in the hallway}
\choiceD{Putting a survey online for anyone to fill out}
\correctAnswer{A}
\explanation{Systematic sampling selects every $n$th member from an ordered list, giving everyone an equal chance.}
\end{practiceQuestion}

\begin{practiceQuestion}{m-8-1-q16}{mc}
\begin{questionText}
A student designs an experiment to test which fertilizer helps plants grow tallest. She uses the same amount of water, sunlight, and soil for each group. The factor she changes (fertilizer type) is called:
\end{questionText}
\choiceA{The conclusion}
\choiceB{The control}
\choiceC{The variable being tested}
\choiceD{The sample}
\correctAnswer{C}
\explanation{The variable being tested (or independent variable) is the factor you deliberately change — here, the fertilizer type.}
\end{practiceQuestion}

\begin{practiceQuestion}{m-8-1-q17}{mc}
\begin{questionText}
A good experimental plan includes keeping all conditions the same except for the one you are testing. This is important because:
\end{questionText}
\choiceA{It makes the experiment longer}
\choiceB{It ensures that differences in results are due to the variable being tested}
\choiceC{It makes the experiment easier}
\choiceD{It is not actually important}
\correctAnswer{B}
\explanation{Controlling other conditions ensures that any difference you observe is caused by the variable you changed, not by something else.}
\end{practiceQuestion}

\begin{practiceQuestion}{m-8-1-q18}{mc}
\begin{questionText}
A company wants to test whether a new package design increases sales. They use the new design in $10$ randomly selected stores and the old design in $10$ other randomly selected stores. What is the sample?
\end{questionText}
\choiceA{All stores owned by the company}
\choiceB{The $20$ selected stores}
\choiceC{Only the $10$ stores with the new design}
\choiceD{Only the $10$ stores with the old design}
\correctAnswer{B}
\explanation{The sample includes all $20$ stores used in the experiment to compare the two designs.}
\end{practiceQuestion}

% --- Short Answer Questions (7) ---

\begin{practiceQuestion}{m-8-1-q19}{sa}
\begin{questionText}
List the five steps of designing an experiment in order.
\end{questionText}
\correctAnswer{1) Ask a question, 2) Identify the population, 3) Choose a sampling method, 4) Collect data, 5) Analyze and conclude}
\explanation{These five steps form a structured approach to collecting and interpreting data about a population.}
\end{practiceQuestion}

\begin{practiceQuestion}{m-8-1-q20}{sa}
\begin{questionText}
A school has $1{,}200$ students. A random sample of $60$ found that $24$ play a musical instrument. Predict how many students in the school play an instrument.
\end{questionText}
\correctAnswer{$480$}
\explanation{$\frac{24}{60} = 40\%$. Then $40\%$ of $1{,}200 = 480$.}
\end{practiceQuestion}

\begin{practiceQuestion}{m-8-1-q21}{sa}
\begin{questionText}
Design a fair experiment to find out which type of exercise seventh graders prefer. Write the question and describe how you would choose a sample.
\end{questionText}
\correctAnswer{Question: ``Which type of exercise do seventh graders prefer?'' Sample: Randomly select $50$ seventh graders from the school roster using a random number generator.}
\explanation{A clear question and a random sampling method ensure the experiment is unbiased and representative.}
\end{practiceQuestion}

\begin{practiceQuestion}{m-8-1-q22}{sa}
\begin{questionText}
Explain why a voluntary response sample (e.g., an online poll) is usually biased.
\end{questionText}
\correctAnswer{People with strong opinions are more likely to respond, so the results do not represent the whole population}
\explanation{Voluntary samples attract people who feel strongly, leaving out those with moderate or no opinions.}
\end{practiceQuestion}

\begin{practiceQuestion}{m-8-1-q23}{sa}
\begin{questionText}
A town has $8{,}000$ residents. A random sample of $200$ found that $70$ support building a new pool. Predict the number of supporters in the town.
\end{questionText}
\correctAnswer{$2{,}800$}
\explanation{$\frac{70}{200} = 35\%$. Then $35\%$ of $8{,}000 = 2{,}800$.}
\end{practiceQuestion}

\begin{practiceQuestion}{m-8-1-q24}{sa}
\begin{questionText}
A student surveys $20$ friends to find out the most popular video game. Explain why this is not a good sample of all seventh graders.
\end{questionText}
\correctAnswer{Friends are likely to have similar interests, so this convenience sample does not represent all seventh graders}
\explanation{A sample of friends is not random. Friends tend to share tastes, leading to biased results.}
\end{practiceQuestion}

\begin{practiceQuestion}{m-8-1-q25}{sa}
\begin{questionText}
What is the difference between a population and a sample? Give an example.
\end{questionText}
\correctAnswer{A population is the entire group; a sample is a smaller part. Example: all students (population) vs.\ $50$ surveyed students (sample).}
\explanation{The population is everyone you want to study. The sample is the subset you actually collect data from.}
\end{practiceQuestion}

% --- Graphical Questions (3 GMC + 2 GSA) ---

\begin{practiceQuestion}{m-8-1-q26}{gmc}
\begin{questionText}
The flowchart shows the five steps of designing an experiment. One step is missing.

\begin{center}
\begin{tikzpicture}[node distance=0.8cm,scale=0.8,every node/.style={transform shape}]
  \node[draw,rounded corners,fill=funBlue!15,minimum width=3cm,minimum height=0.7cm] (s1) {1. Ask a question};
  \node[draw,rounded corners,fill=funGreen!15,minimum width=3cm,minimum height=0.7cm,below=of s1] (s2) {2. Identify the population};
  \node[draw,rounded corners,fill=funOrange!15,minimum width=3cm,minimum height=0.7cm,below=of s2] (s3) {3. ???};
  \node[draw,rounded corners,fill=funRed!15,minimum width=3cm,minimum height=0.7cm,below=of s3] (s4) {4. Collect data};
  \node[draw,rounded corners,fill=gray!15,minimum width=3cm,minimum height=0.7cm,below=of s4] (s5) {5. Analyze and conclude};
  \draw[->,thick] (s1) -- (s2);
  \draw[->,thick] (s2) -- (s3);
  \draw[->,thick] (s3) -- (s4);
  \draw[->,thick] (s4) -- (s5);
\end{tikzpicture}
\end{center}

What is the missing Step 3?
\end{questionText}
\choiceA{Make a graph}
\choiceB{Choose a sampling method}
\choiceC{Write a report}
\choiceD{Find the mean}
\correctAnswer{B}
\explanation{After identifying the population, you must decide HOW to select your sample (randomly, systematically, etc.) before collecting data.}
\end{practiceQuestion}

\begin{practiceQuestion}{m-8-1-q27}{gmc}
\begin{questionText}
The table summarizes four different surveys about student preferences for school lunch.

\begin{center}
\begin{tabular}{|c|l|c|c|}
\hline
\textbf{Survey} & \textbf{Method} & \textbf{Sample Size} & \textbf{\% Pizza} \\
\hline
A & Random from roster & $100$ & $45\%$ \\
\hline
B & Volunteers only & $50$ & $70\%$ \\
\hline
C & One math class & $30$ & $40\%$ \\
\hline
D & Random from roster & $100$ & $48\%$ \\
\hline
\end{tabular}
\end{center}

Which surveys are most reliable for predicting the whole school's preferences?
\end{questionText}
\choiceA{Surveys B and C}
\choiceB{Surveys A and D}
\choiceC{Only Survey B because pizza got the highest percent}
\choiceD{All surveys are equally reliable}
\correctAnswer{B}
\explanation{Surveys A and D use random samples from the full roster and have larger sizes, making them the most representative and reliable.}
\end{practiceQuestion}

\begin{practiceQuestion}{m-8-1-q28}{gmc}
\begin{questionText}
The bar graph shows the results of randomly surveying $50$ students about their favorite subject.

\begin{center}
\begin{tikzpicture}[scale=0.6]
  \draw[thick,->] (0,0) -- (0,5.5) node[above]{\small Students};
  \draw[thick,->] (0,0) -- (9,0);
  \fill[funBlue!70] (0.5,0) rectangle (1.5,4);
  \fill[funGreen!70] (2.5,0) rectangle (3.5,3);
  \fill[funOrange!70] (4.5,0) rectangle (5.5,2);
  \fill[funRed!70] (6.5,0) rectangle (7.5,1);
  \node[below,font=\small] at (1,0) {Math};
  \node[below,font=\small] at (3,0) {Science};
  \node[below,font=\small] at (5,0) {English};
  \node[below,font=\small] at (7,0) {Art};
  \foreach \y in {1,...,5} { \draw (-0.1,\y) -- (0.1,\y); \node[left,font=\small] at (0,\y) {$\y\times 4$}; }
  \node[above,font=\small] at (1,4) {$16$};
  \node[above,font=\small] at (3,3) {$12$};
  \node[above,font=\small] at (5,2) {$8$};
  \node[above,font=\small] at (7,1) {$4$};
\end{tikzpicture}
\end{center}

If the school has $750$ students, predict how many prefer science. Note: $10$ students in the sample did not name one of these four subjects.
\end{questionText}
\choiceA{$12$}
\choiceB{$120$}
\choiceC{$150$}
\choiceD{$180$}
\correctAnswer{D}
\explanation{$\frac{12}{50} = 24\%$. Then $24\%$ of $750 = 180$.}
\end{practiceQuestion}

\begin{practiceQuestion}{m-8-1-q29}{gsa}
\begin{questionText}
A student designed an experiment to find out which snack seventh graders prefer. Read the plan below and identify one problem.

\begin{center}
\begin{tabular}{|l|p{8cm}|}
\hline
\textbf{Question} & Which snack do seventh graders prefer? \\
\hline
\textbf{Population} & All seventh graders at the school \\
\hline
\textbf{Method} & Ask the $15$ students on the basketball team \\
\hline
\textbf{Data Collection} & Record each student's answer \\
\hline
\textbf{Analysis} & Tally results and find the most popular snack \\
\hline
\end{tabular}
\end{center}

Identify the problem with this plan and suggest a better sampling method.
\end{questionText}
\correctAnswer{The sample is biased — basketball players may not represent all seventh graders. A better method is to randomly select $30$--$50$ students from the full class list.}
\explanation{The basketball team is a specific group with potentially different snack preferences (e.g., more protein bars). A random sample from all seventh graders avoids this bias.}
\end{practiceQuestion}

\begin{practiceQuestion}{m-8-1-q30}{gsa}
\begin{questionText}
The table shows the results of a random sample of $40$ students asked about homework time per night.

\begin{center}
\begin{tabular}{|c|c|}
\hline
\textbf{Hours} & \textbf{Students} \\
\hline
$0$--$0.5$ & $8$ \\
\hline
$0.5$--$1$ & $14$ \\
\hline
$1$--$1.5$ & $10$ \\
\hline
$1.5$--$2$ & $6$ \\
\hline
More than $2$ & $2$ \\
\hline
\end{tabular}
\end{center}

The school has $800$ students. Predict how many students spend more than $1$ hour on homework per night.
\end{questionText}
\correctAnswer{$360$ students}
\explanation{Students spending more than $1$ hour: $10 + 6 + 2 = 18$ out of $40 = 45\%$. Then $45\%$ of $800 = 360$.}
\end{practiceQuestion}
